Il lavoro ha presentato Graph Dynamic Meshing come pipeline di $r$-adattività neurale a topologia fissa per una famiglia di PDE transitorie di convezione--diffusione. La pipeline combina una GNN che propone spostamenti nodali dallo stato FEM corrente, un filtro geometrico deterministico lungo il percorso e un solutore ALE--FEM nativo. Il contributo non viene identificato nella sola idea di muovere una mesh con una rete, già affrontata dalla letteratura, ma nella verifica integrata di accuratezza PDE, ammissibilità geometrica e compatibilità ALE nel protocollo considerato.

La campagna E1 ha fornito evidenza di ordine temporale prossimo a uno, di ordini spaziali compatibili con $P_1$ sulla coppia $16^2$--$32^2$ e di consistenza D-GCL in un test indipendente. La convergenza spaziale al livello $64^2$ non è ancora isolata dall'errore temporale e il relativo ordine osservato non viene considerato conclusivo. La SUPG, pur disponibile nel core come opzione, non ha migliorato l'errore nella configurazione MMS testata.

Nella campagna F1, su 28 casi held-out e a parità di gradi di libertà e discretizzazione, GDM ha ottenuto un errore relativo $L^2$ integrato medio del 9.831\%, contro 10.633\% per Fixed FEM e 10.358--10.578\% per i tre controllori classici. La riduzione relativa media appaiata è stata 8.04\% rispetto alla mesh fissa e 6.10--7.81\% rispetto alle baseline classiche, con segno favorevole a GDM in tutti i 28 casi. Questi risultati sono promettenti, ma restano preliminari rispetto alla chiusura dei controlli: alcune baseline hanno raggiunto il limite iterativo prefissato e il gate finale della campagna ha rilevato scarti di arrotondamento rispetto alla soglia geometrica assoluta.

Le diagnostiche validation mostrano che la mesh si concentra moderatamente intorno alla struttura gaussiana mobile, con arricchimento mediano del 21.3\% nella corona di interesse. Il filtro interviene frequentemente e conserva circa il 48.1\% della norma aggregata degli spostamenti grezzi. GDM presenta inoltre qualità minima degli elementi inferiore alle baseline e un costo online circa 12.6 volte quello della mesh fissa. Non viene quindi rivendicato un vantaggio computazionale generale; il risultato principale è un miglioramento di accuratezza nel protocollo testato, ottenuto con un compromesso geometrico e temporale che deve essere ulteriormente caratterizzato.

La validità delle conclusioni è limitata alla famiglia MMS, al dominio e alla risoluzione considerati. Prima della sottomissione sono necessari il completamento della separazione spazio--tempo, l'audit finale del filtro, la convergenza delle baseline classiche e la pubblicazione degli artefatti riproducibili. Il lavoro costituisce pertanto una valutazione integrata e circoscritta di GDM, non una dimostrazione di superiorità universale delle reti neurali per l'adattività di mesh.
